\chapter{Methodology}

The ideal solution would evaluate managed coding agents on a curated benchmark corpus with explicit ground-truth tasks, independent reviewers, multiple agent vendors, repeated runs, preserved prompts, and statistically justified comparisons. It would separate repository health from agent quality, because a failing build in a target repository is not automatically an agent failure. It would also distinguish static architecture classification, dynamic test execution, runtime telemetry, and coverage evidence.

The local constraints are different. The experiment is private, AI-only, and constrained to existing local repositories and configured tools. The method therefore uses a staged evidence pipeline. First, each corpus target is described in JSON with path, kind, expected build and test behavior, coverage requirement, telemetry requirement, and ground-truth status. Second, the runner captures the live git commit and dirty state before executing commands. Third, the runner executes restore, build, and test steps using the repository's native .NET or Node tooling. Fourth, coverage is attempted with the XPlat Code Coverage collector for .NET targets. Fifth, semantic-convention verification is executed separately for qyl's OpenTelemetry packages and the local semconv testbed. Sixth, results are aggregated into JSON and rendered into thesis tables.

qyl and Arqio were selected because they match the research question. qyl contains the automation, OpenTelemetry, dashboard, collector, and semantic-convention surfaces needed to evaluate instrumentation-aware research infrastructure. Arqio contains architecture-analysis assets, benchmark fixtures, metrics files, and dependency-analysis code. Alternatives were considered implicitly by the available local workspace: a standalone thesis-only script would be simpler, but it would not exercise the same automation and observability surfaces; a cloud benchmark would be broader, but it would introduce credentials and non-local reproducibility risk.

The evaluation method uses four quality dimensions. Build and test status indicate whether each target can be exercised by the harness. Coverage status indicates whether test execution produced measurable code-coverage evidence. Semantic-convention status checks whether qyl's telemetry naming surface can be verified against local semconv tooling without publishing packages. Reproducibility is estimated through repeated target runs when time permits; when only one run is available, the result is explicitly marked as preliminary. Classification output is reported as agreement, not accuracy, unless a target has a configured ground-truth label.
